\NeedsTeXFormat{LaTeX2e}[1995/12/01]
\ProvidesClass{sduthesis}[2024/12/01 1.1.0 Shandong University Thesis Template]

%% \iffalse
%%
%% ===================================================================
%% SDUTeX - 山东大学 LaTeX 论文模板核心包
%% https://github.com/h1s97x/sdutex
%% ===================================================================
%%
%% Copyright (C) 2024 SDU LaTeX Contributors
%%
%% This work may be distributed and/or modified under the conditions
%% of the LaTeX Project Public License, either version 1.3 of this
%% license or (at your option) any later version.
%% The latest version of this license is in
%%   https://www.latex-project.org/lppl.txt
%% and version 1.3c or later is part of all distributions of LaTeX
%% version 2008-05-04 or later.
%%
%% This work has the LPPL maintenance status "maintained".
%%
%% The Current Maintainer of this work is SDU LaTeX Contributors.
%%
%% This work consists of the files listed in the companion file
%% sduthesis.dtx and the other files listed in the header of that file.
%%
%% The source files of this work are located at:
%%   https://github.com/h1s97x/sdutex
%%
%% \fi
%%
%% ===================================================================
%% Class Structure
%% ===================================================================
%%
%% This document class implements the *kernel* of the SDUTeX plugin
%% architecture:
%%
%%   1. Engine + base class loading + basic layout
%%   2. The SDUTeX Hook system (6 named hooks)
%%   3. Module loader (``module = {...}'') triggered at begindocument
%%
%% All degree-specific logic (cover pages, blind-review hiding, etc.)
%% has been moved out into the companion ``modules/*.sty'' files so that
%% the kernel stays degree-agnostic and acts as the single source of
%% truth.
%%
%% \subsection{Class Identity}
%%
%% \changes{v1.1.0}{2024/12/01}{Kernel + modules + hooks architecture}
%%
%%    \begin{macrocode}
%<!CLASS>
%    \end{macrocode}
%%
%%    \begin{macrocode}
%<!CLASS>
\RequirePackage{expl3}
\RequirePackage{l3keys2e}
\ExplSyntaxOn
%    \end{macrocode}
%%
%% Define the internal prefix.
%%    \begin{macrocode}
%<!CLASS>
%<@@=sdu>
%    \end{macrocode}
%%
%% ===================================================================
%% Global Variables
%% ===================================================================
%%
%% \subsection{Global Variables}
%%
%% Degree type (used only as a convenience shortcut that maps onto the
%% module loader):
%% \begin{itemize}
%%   \item 1 = bachelor (本科生)
%%   \item 2 = master (学术型硕士)
%%   \item 3 = master-professional (专业型硕士)
%%   \item 4 = doctor (博士)
%% \end{itemize}
%%    \begin{macrocode}
%<!CLASS>
\int_new:N \g_@@_type_int
\int_new:N \g_@@_type_major_int
\int_new:N \g_@@_type_professional_int
%    \end{macrocode}
%%
%% \begin{variable}{\g_@@_blind_bool}
%% 是否为盲审模式
%%    \begin{macrocode}
%<!CLASS>
\bool_new:N \g_@@_blind_bool
%    \end{macrocode}
%% \end{variable}
%%
%% \begin{variable}{\g_@@_twoside_bool}
%% 是否双面打印
%%    \begin{macrocode}
%<!CLASS>
\bool_new:N \g_@@_twoside_bool
%    \end{macrocode}
%% \end{variable}
%%
%% \begin{variable}{\g_@@_professional_bool}
%% 是否为专业学位（由 master 模块使用）
%%    \begin{macrocode}
%<!CLASS>
\bool_new:N \g_@@_professional_bool
%    \end{macrocode}
%% \end{variable}
%%
%% \begin{variable}{\g_@@_module_clist}
%% 需要加载的模块列表（如 master, blindreview）
%%    \begin{macrocode}
%<!CLASS>
\clist_new:N \g_@@_module_clist
\clist_new:N \g_@@_module_pending_clist
%    \end{macrocode}
%% \end{variable}
%%
%% User configuration variables.
%%    \begin{macrocode}
%<!CLASS>
\tl_new:N \g_@@_title_tl
\tl_new:N \g_@@_title_en_tl
\tl_new:N \g_@@_author_tl
\tl_new:N \g_@@_author_en_tl
\tl_new:N \g_@@_student_id_tl
\tl_new:N \g_@@_school_tl
\tl_new:N \g_@@_major_tl
\tl_new:N \g_@@_major_en_tl
\tl_new:N \g_@@_supervisor_tl
\tl_new:N \g_@@_supervisor_title_tl
\tl_new:N \g_@@_co_supervisor_tl
\tl_new:N \g_@@_supervisor_en_tl
\tl_new:N \g_@@_co_supervisor_en_tl
\tl_new:N \g_@@_date_tl
\tl_new:N \g_@@_cn_keywords_tl
\tl_new:N \g_@@_en_keywords_tl
\tl_new:N \g_@@_abstract_tl
\tl_new:N \g_@@_abstract_en_tl
\tl_new:N \g_@@_degree_name_tl
\tl_new:N \g_@@_degree_name_en_tl
% 扩展选项变量
\tl_new:N \g_@@_lang_tl
\tl_new:N \g_@@_math_style_tl
\tl_new:N \g_@@_preset_tl
\bool_new:N \g_@@_has_listoffigures_bool
\bool_new:N \g_@@_has_listoftables_bool
\tl_new:N \g_@@_header_left_tl
\tl_new:N \g_@@_header_right_tl
%    \end{macrocode}
%%
%% Layout variables.
%%    \begin{macrocode}
%<!CLASS>
\dim_new:N \g_@@_cover_label_width_dim
\dim_new:N \g_@@_cover_value_width_dim
%    \end{macrocode}
%%
%% ===================================================================
%% Auxiliary Functions
%% ===================================================================
%%
%% \subsection{Auxiliary Functions}
%%
%% The kernel keeps only the degree-agnostic helpers.  Degree-specific
%% predicates (bachelor / master / doctor / professional) are provided
%% by the degree modules so that the kernel itself stays neutral.
%%
%% \begin{function}{\@@_if_blind:TF}
%% 判断是否为盲审模式
%%    \begin{macrocode}
%<!CLASS>
\cs_new:Npn \@@_if_blind:TF #1#2 {
  \bool_if:NTF \g_@@_blind_bool {#1} {#2}
}
%    \end{macrocode}
%% \end{function}
%%
%% \begin{function}{\@@_secret_info:n}
%% 隐藏敏感信息（用于盲审），由 blindreview 模块的钩子改写。
%%    \begin{macrocode}
%<!CLASS>
\cs_new:Npn \@@_secret_circle:n #1 {
  \leavevmode \prg_replicate:nn { #1 } {
    \raisebox{0.2ex}{\rule{0.9em}{0.9em}}\kern 0.1em
  }
}
\cs_new:Npn \@@_secret_info:n #1 {
  \@@_if_blind:TF {
    \@@_secret_circle:n { 1 }
  } {
    #1
  }
}
%    \end{macrocode}
%% \end{function}
%%
%% \begin{function}{\@@_secret_info_fixed:Nn}
%% 固定长度隐藏（用于学号等）
%%    \begin{macrocode}
%<!CLASS>
\cs_new:Npn \@@_secret_info_fixed:Nn #1#2 {
  \tl_if_blank:VTF #1 { } {
    \@@_if_blind:TF {
      \@@_secret_circle:n {#2}
    } {
      #1
    }
  }
}
%    \end{macrocode}
%% \end{function}
%%
%% ===================================================================
%% Cover Helper
%% ===================================================================
%%
%% \subsection{Cover Helper}
%%
%% A shared, degree-agnostic cover table entry used by the degree
%% modules.  This is basic layout infrastructure and therefore stays in
%% the kernel.
%%    \begin{macrocode}
%<!CLASS>
\dim_gset:Nn \g_@@_cover_label_width_dim { 2.3cm }
\dim_gset:Nn \g_@@_cover_value_width_dim { 5.5cm }
%    \end{macrocode}
%%
%%    \begin{macrocode}
%<!CLASS>
\cs_new:Npn \@@_cover_entry:nn #1#2 {
  \begin{tabular}{>{\centering\arraybackslash}p{\g_@@_cover_label_width_dim}>{\centering\arraybackslash}p{\g_@@_cover_value_width_dim}}
    \zihao{3} \bfseries \kaishu {#1} &
    \zihao{3} \kaishu {#2}
    \\ \cline{2-2}
  \end{tabular}
  \vspace{-5mm}
}
%    \end{macrocode}
%%
%% ===================================================================
%% Hook System
%% ===================================================================
%%
%% \subsection{Hook System}
%%
%% SDUTeX 插件化架构的核心：定义 6 个命名钩子，供模块与用户扩展。
%% 钩子基于 LaTeX3 内核的 hook 机制实现，支持按标签追加代码。
%%
%% \begin{function}{sduthesis/after-setup}
%% 在 \cs{SDUSetup} 处理完配置后触发。
%%    \begin{macrocode}
%<!CLASS>
\hook_new:n { sduthesis/after-setup }
\hook_new:n { sduthesis/before-cover }
\hook_new:n { sduthesis/cover-style }
\hook_new:n { sduthesis/frontmatter/begin }
\hook_new:n { sduthesis/mainmatter/begin }
\hook_new:n { sduthesis/backmatter/begin }
%    \end{macrocode}
%% \end{function}
%%
%% \begin{function}{\sdu_hook_use:n}
%% 触发一个命名钩子。
%%    \begin{macrocode}
%<!CLASS>
\cs_new:Npn \sdu_hook_use:n #1 {
  \hook_use:n { #1 }
}
%    \end{macrocode}
%% \end{function}
%%
%% ===================================================================
%% Options Parsing
%% ===================================================================
%%
%% \subsection{Options Parsing}
%%
%% Declare class options.  The \texttt{degree} and \texttt{blind}
%% options are kept for backward compatibility and transparently mapped
%% onto the module loader; \texttt{module} provides the new combined
%% loading interface.
%%    \begin{macrocode}
%<!CLASS>
\keys_define:nn { sdu / option } {
  degree .choices:nn = { bachelor, master, professional, doctor } {
    \int_gset:Nn \g_@@_type_int {\l_keys_choice_int}
    \bool_gset_false:N \g_@@_professional_bool
    \str_if_eq:eeT { \l_keys_choice_tl } { professional } {
      \bool_gset_true:N \g_@@_professional_bool
    }
    \clist_gput_right:Nx \g_@@_module_pending_clist { degree-\l_keys_choice_tl }
  },
  degree .value_required:n = true,
  degree .initial:n = {bachelor},
  professional .bool_set:N = \g_@@_professional_bool,
  blind .bool_set:N = \g_@@_blind_bool,
  blind .initial:n = {false},
  twoside .bool_set:N = \g_@@_twoside_bool,
  twoside .initial:n = {false},
  module .code:n = {
    \clist_gput_right:Nn \g_@@_module_pending_clist { #1 }
  },
  % 语言选项
  lang .choices:nn = { zh, en, zh-en, en-zh } {
    \tl_gset:Nx \g_@@_lang_tl { \l_keys_choice_tl }
  },
  lang .initial:n = { zh },
  % 数学风格选项
  math-style .choices:nn = { ISO, GB, french, upright } {
    \tl_gset:Nx \g_@@_math_style_tl { \l_keys_choice_tl }
  },
  math-style .initial:n = { ISO },
  % 预设风格选项
  preset .choices:nn = { sdu, sjtu, none } {
    \tl_gset:Nx \g_@@_preset_tl { \l_keys_choice_tl }
  },
  preset .initial:n = { sdu },
  unknown .code:n = {
    \PassOptionsToClass {\l_keys_key_tl} {ctexbook}
  }
}
%    \end{macrocode}
%%
%% Process options.
%%    \begin{macrocode}
%<!CLASS>
\ProcessKeysOptions { sdu / option }
%    \end{macrocode}
%%
%% Legacy command to set the degree type (maps onto module loading).
%%    \begin{macrocode}
%<!CLASS>
\NewDocumentCommand \degree { m } {
  \keys_set:nn { sdu / option } { degree = {#1} }
}
%    \end{macrocode}
%%
%% ===================================================================
%% Load Base Class
%% ===================================================================
%%
%% \subsection{Load Base Class}
%%
%%    \begin{macrocode}
%<!CLASS>
\LoadClass[zihao=-4, a4paper, twoside=\g_@@_twoside_bool] {ctexbook}
%    \end{macrocode}
%%
%% ===================================================================
%% Package Loading
%% ===================================================================
%%
%% \subsection{Package Loading}
%%
%%    \begin{macrocode}
%<!CLASS>
\RequirePackage { geometry }
\RequirePackage { fancyhdr }
\RequirePackage { graphicx }
\RequirePackage { array }
\RequirePackage { booktabs }
\RequirePackage { longtable }
\RequirePackage { enumitem }
\RequirePackage { caption }
\RequirePackage { hyperref }
\RequirePackage { amsmath }
\RequirePackage { amssymb }
\RequirePackage { amsthm }
\RequirePackage { bm }
\RequirePackage { algorithmicx }
\RequirePackage { algorithm }
\RequirePackage { listings }
\RequirePackage { xcolor }
\RequirePackage { setspace }
% 中文下划线/着重号：XeTeX 用 xeCJKfntef，LuaTeX 用 luatexja 兼容层
\sys_if_engine_xetex:TF {
  \RequirePackage { xeCJKfntef }
} {
  \RequirePackage { luatexja }
}
%    \end{macrocode}
%%
%% ===================================================================
%% Page Layout
%% ===================================================================
%%
%% \subsection{Page Layout}
%%
%% Set page geometry.
%%    \begin{macrocode}
%<!CLASS>
\geometry {
  top = 3.0cm,
  bottom = 2.5cm,
  left = 3.0cm,
  right = 2.5cm,
  headheight = 1.5cm,
  headsep = 0.5cm,
  footskip = 1.0cm
}
%    \end{macrocode}
%%
%% ===================================================================
%% Style Configuration
%% ===================================================================
%%
%% \subsection{Style Configuration}
%%
%% Set Chinese font (using system fonts).
%%    \begin{macrocode}
%<!CLASS>
% Windows 环境优先使用系统字体（Times New Roman / Arial / Consolas），
% 其他平台（Linux CI / macOS）回退到 TeX Live 自带的 TeX Gyre 字体，
% 以保证 make test 在无 Windows 字体的环境中也能通过。
\sys_if_platform_windows:TF {
  \setmainfont{Times~New~Roman}[AutoFakeBold]
  \setsansfont{Arial}
  \setmonofont{Consolas}
} {
  \setmainfont{TeX~Gyre~Termes}[AutoFakeBold]
  \setsansfont{TeX~Gyre~Heros}
  \setmonofont{TeX~Gyre~Cursor}
}
%    \end{macrocode}
%%
%% Set bibliography style (use custom .bst).
%%    \begin{macrocode}
%<!CLASS>
\bibliographystyle{sduthesis}
%    \end{macrocode}
%%
%% Set algorithm style.
%%    \begin{macrocode}
%<!CLASS>
\RequirePackage { algpseudocode }
%    \end{macrocode}
%%
%% ===================================================================
%% Theorem Environments
%% ===================================================================
%%
%% \subsection{Theorem Environments}
%%
%% Define standard theorem-like environments.
%%    \begin{macrocode}
%<!CLASS>
\newtheorem{definition}{定义}
\newtheorem{theorem}{定理}
\newtheorem{lemma}{引理}
\newtheorem{corollary}{推论}
\newtheorem{proposition}{命题}
% 使用 amsthm 内置的 proof 环境
%    \end{macrocode}
%%
%% ===================================================================
%% User Setup Interface
%% ===================================================================
%%
%% \subsection{User Setup Interface}
%%
%% Define setup keys.  \texttt{module} 用于在文档内组合加载模块，
%% 如 \verb|\SDUSetup{ module = {master, blindreview} }|。
%%    \begin{macrocode}
%<!CLASS>
\keys_define:nn { sdu / setup } {
  title .tl_gset:N = \g_@@_title_tl,
  title* .tl_gset:N = \g_@@_title_en_tl,
  author .tl_gset:N = \g_@@_author_tl,
  author* .tl_gset:N = \g_@@_author_en_tl,
  student-id .tl_gset:N = \g_@@_student_id_tl,
  school .tl_gset:N = \g_@@_school_tl,
  major .tl_gset:N = \g_@@_major_tl,
  major* .tl_gset:N = \g_@@_major_en_tl,
  supervisor .tl_gset:N = \g_@@_supervisor_tl,
  supervisor-title .tl_gset:N = \g_@@_supervisor_title_tl,
  supervisor* .tl_gset:N = \g_@@_supervisor_en_tl,
  co-supervisor .tl_gset:N = \g_@@_co_supervisor_tl,
  co-supervisor* .tl_gset:N = \g_@@_co_supervisor_en_tl,
  date .tl_gset:N = \g_@@_date_tl,
  keywords .tl_gset:N = \g_@@_cn_keywords_tl,
  keywords* .tl_gset:N = \g_@@_en_keywords_tl,
  % 双语分组键值（对齐 SJTUTeX）：\verb|info = { zh/title = {...}, en/title = {...} }|
  info .code:n = {
    \keys_set:nn { sdu / info } { #1 }
  },
  % 语言、数学风格、预设
  lang .choices:nn = { zh, en, zh-en, en-zh } {
    \tl_gset:Nx \g_@@_lang_tl { \l_keys_choice_tl }
  },
  math-style .choices:nn = { ISO, GB, french, upright } {
    \tl_gset:Nx \g_@@_math_style_tl { \l_keys_choice_tl }
  },
  preset .choices:nn = { sdu, sjtu, none } {
    \tl_gset:Nx \g_@@_preset_tl { \l_keys_choice_tl }
  },
  % 页眉样式设置
  header-left .tl_gset:N = \g_@@_header_left_tl,
  header-right .tl_gset:N = \g_@@_header_right_tl,
  module .code:n = {
    \clist_gput_right:Nn \g_@@_module_pending_clist { #1 }
  },
}

% 双语分组：\verb|info = { zh/title = {...}, en/title = {...} }|
\keys_define:nn { sdu / info } {
  zh/title .tl_gset:N = \g_@@_title_tl,
  zh/author .tl_gset:N = \g_@@_author_tl,
  zh/student-id .tl_gset:N = \g_@@_student_id_tl,
  zh/school .tl_gset:N = \g_@@_school_tl,
  zh/major .tl_gset:N = \g_@@_major_tl,
  zh/supervisor .tl_gset:N = \g_@@_supervisor_tl,
  zh/supervisor-title .tl_gset:N = \g_@@_supervisor_title_tl,
  zh/co-supervisor .tl_gset:N = \g_@@_co_supervisor_tl,
  zh/date .tl_gset:N = \g_@@_date_tl,
  zh/keywords .tl_gset:N = \g_@@_cn_keywords_tl,
  zh/degree-name .tl_gset:N = \g_@@_degree_name_tl,
  en/title .tl_gset:N = \g_@@_title_en_tl,
  en/author .tl_gset:N = \g_@@_author_en_tl,
  en/major .tl_gset:N = \g_@@_major_en_tl,
  en/supervisor .tl_gset:N = \g_@@_supervisor_en_tl,
  en/co-supervisor .tl_gset:N = \g_@@_co_supervisor_en_tl,
  en/keywords .tl_gset:N = \g_@@_en_keywords_tl,
  en/degree-name .tl_gset:N = \g_@@_degree_name_en_tl,
}
%    \end{macrocode}
%%
%% \begin{function}{\SDUSetup}
%% Main setup command for users.  Fires the
%% \texttt{sduthesis/after-setup} hook once configuration is applied.
%%    \begin{macrocode}
%<!CLASS>
\NewDocumentCommand \SDUSetup { m } {
  \keys_set:nn { sdu / setup } {#1}
  \hook_use:n { sduthesis/after-setup }
}
%    \end{macrocode}
%% \end{function}
%%
%% Legacy commands mapping to the modern setup keys.
%%    \begin{macrocode}
%<!CLASS>
\DeclareDocumentCommand \title { s m } {
  \IfBooleanTF {#1} {
    \keys_set:nn { sdu / setup } { title* = {#2} }
  } {
    \keys_set:nn { sdu / setup } { title = {#2} }
  }
}
\DeclareDocumentCommand \author { s m } {
  \IfBooleanTF {#1} {
    \keys_set:nn { sdu / setup } { author* = {#2} }
  } {
    \keys_set:nn { sdu / setup } { author = {#2} }
  }
}
\DeclareDocumentCommand \school { m } { \keys_set:nn { sdu / setup } { school = {#1} } }
\DeclareDocumentCommand \major { s m } {
  \IfBooleanTF {#1} {
    \keys_set:nn { sdu / setup } { major* = {#2} }
  } {
    \keys_set:nn { sdu / setup } { major = {#2} }
  }
}
\DeclareDocumentCommand \supervisor { s m } {
  \IfBooleanTF {#1} {
    \keys_set:nn { sdu / setup } { supervisor* = {#2} }
  } {
    \keys_set:nn { sdu / setup } { supervisor = {#2} }
  }
}
\DeclareDocumentCommand \cosupervisor { s m } {
  \IfBooleanTF {#1} {
    \keys_set:nn { sdu / setup } { co-supervisor* = {#2} }
  } {
    \keys_set:nn { sdu / setup } { co-supervisor = {#2} }
  }
}
\DeclareDocumentCommand \date { m } { \keys_set:nn { sdu / setup } { date = {#1} } }
\DeclareDocumentCommand \studentid { m } { \keys_set:nn { sdu / setup } { student-id = {#1} } }
\DeclareDocumentCommand \keywords { s m } {
  \IfBooleanTF {#1} {
    \keys_set:nn { sdu / setup } { keywords* = {#2} }
  } {
    \keys_set:nn { sdu / setup } { keywords = {#2} }
  }
}
%    \end{macrocode}
%%
%% Getter commands.
%%    \begin{macrocode}
%<!CLASS>
\NewDocumentCommand \GetTitle { } { \tl_use:N \g_@@_title_tl }
\NewDocumentCommand \GetTitleEn { } { \tl_use:N \g_@@_title_en_tl }
\NewDocumentCommand \GetAuthor { } { \tl_use:N \g_@@_author_tl }
\NewDocumentCommand \GetAuthorEn { } { \tl_use:N \g_@@_author_en_tl }
\NewDocumentCommand \GetStudentId { } { \tl_use:N \g_@@_student_id_tl }
\NewDocumentCommand \GetSchool { } { \tl_use:N \g_@@_school_tl }
\NewDocumentCommand \GetMajor { } { \tl_use:N \g_@@_major_tl }
\NewDocumentCommand \GetMajorEn { } { \tl_use:N \g_@@_major_en_tl }
\NewDocumentCommand \GetSupervisor { } { \tl_use:N \g_@@_supervisor_tl }
\NewDocumentCommand \GetSupervisorTitle { } { \tl_use:N \g_@@_supervisor_title_tl }
\NewDocumentCommand \GetSupervisorEn { } { \tl_use:N \g_@@_supervisor_en_tl }
\NewDocumentCommand \GetCoSupervisorEn { } { \tl_use:N \g_@@_co_supervisor_en_tl }
\NewDocumentCommand \GetCoSupervisor { } { \tl_use:N \g_@@_co_supervisor_tl }
\NewDocumentCommand \GetDate { } { \tl_use:N \g_@@_date_tl }
\NewDocumentCommand \GetCnKeywords { } { \tl_use:N \g_@@_cn_keywords_tl }
\NewDocumentCommand \GetEnKeywords { } { \tl_use:N \g_@@_en_keywords_tl }
% 新增 Getter：语言 / 数学风格 / 预设 / 学位名称
\NewDocumentCommand \GetLang { } { \tl_use:N \g_@@_lang_tl }
\NewDocumentCommand \GetMathStyle { } { \tl_use:N \g_@@_math_style_tl }
\NewDocumentCommand \GetPreset { } { \tl_use:N \g_@@_preset_tl }
\NewDocumentCommand \GetDegreeName { } { \tl_use:N \g_@@_degree_name_tl }
\NewDocumentCommand \GetDegreeNameEn { } { \tl_use:N \g_@@_degree_name_en_tl }
%    \end{macrocode}
%%
%% ===================================================================
%% Module Loader
%% ===================================================================
%%
%% \subsection{Module Loader}
%%
%% Load a single module file \texttt{sduthesis-<name>.sty}.
%%    \begin{macrocode}
%<!CLASS>
\cs_new_protected:Npn \sdu_load_module:n #1 {
  \file_if_exist:nTF { sduthesis-#1.sty } {
    \RequirePackage { sduthesis-#1 }
    \clist_gput_right:Nx \g_@@_module_clist { #1 }
  } {
    \msg_warning:nnx { sdutex } { unknown-module } { #1 }
  }
}
%    \end{macrocode}
%%
%% \begin{function}{\sdu_load_module:}
%% Public user-level command to load a module at any point.
%%    \begin{macrocode}
%<!CLASS>
\NewDocumentCommand \sdu_load_module: { m } {
  \sdu_load_module:n { #1 }
}
%    \end{macrocode}
%% \end{function}
%%
%% Resolve the pending module requests into a deduplicated load list and
%% load each module.  Runs inside \cs{AtBeginDocument} so that module
%% files can rely on the class being fully available.
%%
%% \begin{function}{\@@_resolve_module:n}
%% Convert a raw pending entry (either a real module name or a
%% \texttt{degree-<name>} placeholder) into the module name(s) to load.
%%    \begin{macrocode}
%<!CLASS>
\cs_new_protected:Npn \@@_resolve_module:n #1 {
  \str_case:nn { #1 } {
    { degree-bachelor } { \clist_gput_right:Nn \g_@@_module_clist { undergraduate } }
    { degree-master } { \clist_gput_right:Nn \g_@@_module_clist { master } }
    { degree-professional } { \clist_gput_right:Nn \g_@@_module_clist { master } }
    { degree-doctor } { \clist_gput_right:Nn \g_@@_module_clist { doctor } }
  } {
    \clist_gput_right:Nn \g_@@_module_clist { #1 }
  }
}
%    \end{macrocode}
%% \end{function}
%%
%% \begin{function}{\@@_load_pending_modules:}
%%    \begin{macrocode}
%<!CLASS>
\cs_new_protected:Npn \@@_load_pending_modules: {
  \clist_map_inline:Nn \g_@@_module_pending_clist {
    \@@_resolve_module:n { ##1 }
  }
  \clist_clear:N \g_@@_module_pending_clist
  % 去重后加载
  \clist_map_inline:Nn \g_@@_module_clist {
    \sdu_load_module:n { ##1 }
  }
  \clist_clear:N \g_@@_module_clist
}
%    \end{macrocode}
%% \end{function}
%%
%%    \begin{macrocode}
%<!CLASS>
\AtBeginDocument {
  \@@_load_pending_modules:
}
%    \end{macrocode}
%%
%% ===================================================================
%% Front/back matter hooks
%% ===================================================================
%%
%% \subsection{Front/back matter hooks}
%%
%% The kernel exposes \cs{frontmatter}, \cs{mainmatter} and
%% \cs{backmatter} so the corresponding hooks fire at the right time.
%%    \begin{macrocode}
%<!CLASS>
\RenewDocumentCommand \frontmatter { } {
  \hook_use:n { sduthesis/frontmatter/begin }
  \relax
}
\RenewDocumentCommand \mainmatter { } {
  \hook_use:n { sduthesis/mainmatter/begin }
  \relax
}
\RenewDocumentCommand \backmatter { } {
  \hook_use:n { sduthesis/backmatter/begin }
  \relax
}
%    \end{macrocode}
%%
%% ===================================================================
%% Cover Page Dispatcher
%% ===================================================================
%%
%% \subsection{Cover Page Dispatcher}
%%
%% \cs{MakeCover} triggers the \texttt{before-cover} and
%% \texttt{cover-style} hooks.  The actual cover layout is provided by
%% the loaded degree module via the \texttt{cover-style} hook.
%%    \begin{macrocode}
%<!CLASS>
\NewDocumentCommand \MakeCover { } {
  \hook_use:n { sduthesis/before-cover }
  \hook_use:n { sduthesis/cover-style }
}
%    \end{macrocode}
%%
%% Legacy lowercase alias for \MakeCover.
%%    \begin{macrocode}
%<!CLASS>
\NewDocumentCommand \makecover { } { \MakeCover }
%    \end{macrocode}
%%
%% ===================================================================
%% Abstract Environments
%% ===================================================================
%%
%% \subsection{Abstract Environments}
%%
%% These are shared by every degree type, so they stay in the kernel.
%%
%% Chinese abstract.  支持可选参数设置标题（默认“摘要”）。
%%    \begin{macrocode}
%<!CLASS>
\NewDocumentEnvironment { cnabstract } { o } {
  \phantomsection
  \pdfbookmark[0] { 摘要 } { cnabstract }
  \begin{center}
    \IfValueTF{#1} {
      \zihao{3} \bfseries \heiti { #1 }
    } {
      \zihao{3} \bfseries \heiti { 摘\qquad 要 }
    }
  \end{center}
  \vspace{12pt}
  \setlength{\parindent}{1.2em}
  \zihao{-4}
}{
  \vspace{1cm}
  \par \noindent
  \zihao{-4} \bfseries \heiti { 关键词： }
  \zihao{-4} \GetCnKeywords
  \clearpage
}
%    \end{macrocode}
%%
%% English abstract.  支持可选参数设置标题（默认 ABSTRACT）。
%%    \begin{macrocode}
%<!CLASS>
\NewDocumentEnvironment { enabstract } { o } {
  \phantomsection
  \pdfbookmark[0] { Abstract } { enabstract }
  \begin{center}
    \IfValueTF{#1} {
      \zihao{-2} \bfseries { #1 }
    } {
      \zihao{-2} \bfseries { ABSTRACT }
    }
  \end{center}
  \vspace{12pt}
  \setlength{\parindent}{1.2em}
  \zihao{-4}
}{
  \vspace{1cm}
  \par \noindent
  \zihao{-4} \bfseries { Keywords: }
  \zihao{-4} \GetEnKeywords
  \clearpage
}
%    \end{macrocode}
%%
%% ===================================================================
%% Declaration
%% ===================================================================
%%
%% \subsection{Declaration}
%%
%% Academic declaration (shared by all degrees).
%%    \begin{macrocode}
%<!CLASS>
\NewDocumentCommand \MakeDeclaration { } {
  \phantomsection
  \pdfbookmark[0] { 声明 } { declaration }
  \begin{center}
    \zihao{3} \bfseries \heiti { 学位论文独创性声明 }
  \end{center}
  \vspace{12pt}
  \setlength{\parindent}{2em}
  \zihao{-4}

  本人郑重声明，所呈交的学位论文，是本人在导师指导下，独立进
  行研究工作所取得的研究成果。除文中已经标明引用的内容外，本
  论文不包含任何其他个人或集体已经发表或撰写过的研究成果。对
  本文的研究做出贡献的个人和集体，均已在文中以明确方式标明。

  本声明的法律结果由本人承担。

  \vspace{2cm}
  \hfill
  \begin{tabular}{lcl}
    学位论文作者签名： & \hspace{3cm} & \\
    & & \\
    日期： & \hspace{3cm} & 年 \hspace{1cm} 月 \hspace{1cm} 日 \\
  \end{tabular}

  \clearpage
}
%    \end{macrocode}
%%
%% Legacy lowercase alias for \MakeDeclaration.
%%    \begin{macrocode}
%<!CLASS>
\NewDocumentCommand \makedeclaration { } { \MakeDeclaration }
%    \end{macrocode}
%%
%% ===================================================================
%% Acknowledgement
%% ===================================================================
%%
%% \subsection{Acknowledgement}
%%
%% Acknowledgement environment (shared by all degrees).
%%    \begin{macrocode}
%<!CLASS>
\NewDocumentEnvironment { acknowledgement } { } {
  \phantomsection
  \pdfbookmark[0] { 致谢 } { acknowledgement }
  \begin{center}
    \zihao{3} \bfseries \heiti { 致\qquad 谢 }
  \end{center}
  \vspace{12pt}
  \setlength{\parindent}{2em}
  \zihao{-4}
}{
  \clearpage
}
%    \end{macrocode}
%%
%% ===================================================================
%% Appendix
%% ===================================================================
%%
%% \subsection{Appendix}
%%
%% Appendix command (shared by all degrees).
%%    \begin{macrocode}
%<!CLASS>
\NewDocumentCommand \MakeAppendix { } {
  \appendix
  \counterwithin { equation } { chapter }
  \counterwithin { figure } { chapter }
  \counterwithin { table } { chapter }
}
%    \end{macrocode}
%%
%% ===================================================================
%% Section Style
%% ===================================================================
%%
%% \subsection{Section Style}
%%
%% Configure section headings.
%%    \begin{macrocode}
%<!CLASS>
\ctexset {
  chapter = {
    format = { \heiti \bfseries \centering \zihao{3} },
    beforeskip = { 0pt },
    afterskip = { 18pt },
    number = { \thechapter }
  },
  section = {
    format = { \heiti \bfseries \zihao{4} },
    beforeskip = { 18pt },
    afterskip = { 12pt }
  },
  subsection = {
    format = { \kaishu \bfseries \zihao{-4} },
    beforeskip = { 12pt },
    afterskip = { 6pt }
  }
}
%    \end{macrocode}
%%
%% ===================================================================
%% Header and Footer
%% ===================================================================
%%
%% \subsection{Header and Footer}
%%
%% Configure header and footer.  支持通过
%% \verb|\SDUSetup{ header-left = {...}, header-right = {...} }| 自定义页眉。
%% 默认页眉左侧为论文标题，右侧为章节名。
%%    \begin{macrocode}
%<!CLASS>
\fancypagestyle{plain}{
  \fancyhf{}
  \fancyfoot[C]{\thepage}
  \renewcommand{\headrulewidth}{0pt}
}

\fancypagestyle{sduthesis}{
  \fancyhf{}
  \fancyfoot[C]{\thepage}
  \renewcommand{\headrulewidth}{0.4pt}
  \tl_if_empty:NTF \g_@@_header_left_tl {
    \fancyhead[L]{\zihao{-5} \GetTitle}
  } {
    \fancyhead[L]{\zihao{-5} \tl_use:N \g_@@_header_left_tl}
  }
  \tl_if_empty:NTF \g_@@_header_right_tl {
    \fancyhead[R]{\zihao{-5} \leftmark}
  } {
    \fancyhead[R]{\zihao{-5} \tl_use:N \g_@@_header_right_tl}
  }
}

\pagestyle{plain}
%    \end{macrocode}
%%
%% \subsection{List of Figures and Tables}
%%
%% 图表目录命令（硬缺失功能补充）。
%%    \begin{macrocode}
%<!CLASS>
\NewDocumentCommand \ListOfFigures { } {
  \phantomsection
  \pdfbookmark[0] { 图目录 } { listoffigures }
  \listoffigures
}
\NewDocumentCommand \ListOfTables { } {
  \phantomsection
  \pdfbookmark[0] { 表目录 } { listoftables }
  \listoftables
}
%    \end{macrocode}
%%
%% \subsection{Theorem Environment Enhancement}
%%
%% 增强定理环境：补充常用定理样式，并支持 \verb|math-style| 选项。
%%    \begin{macrocode}
%<!CLASS>
% 根据 math-style 调整数学风格
\str_case:Vn \g_@@_math_style_tl {
  { french } {
    \AtBeginDocument { \let\div\relax }
  }
}
% 定理环境（已在前面定义 definition/theorem/lemma/corollary/proposition）
% 补充：assumption（假设）、conjecture（猜想）、remark（注）、example（例）
\theoremstyle{plain}
\newtheorem{assumption}{假设}
\newtheorem{conjecture}{猜想}
\newtheorem*{remark}{注}
\theoremstyle{definition}
\newtheorem{example}{例}
% 定理环境的章节编号
\counterwithin{theorem}{chapter}
\counterwithin{definition}{chapter}
\counterwithin{lemma}{chapter}
\counterwithin{corollary}{chapter}
\counterwithin{proposition}{chapter}
\counterwithin{assumption}{chapter}
\counterwithin{conjecture}{chapter}
\counterwithin{example}{chapter}
%%
%% Message definitions.
%%    \begin{macrocode}
%<!CLASS>
\msg_new:nnn { sdutex } { unknown-module } {
  Module~'#1'~not~found.~Expected~'sduthesis-#1.sty'.~
  Available~modules:~undergraduate,~master,~doctor,~blindreview.
}
%    \end{macrocode}
%%
%% ===================================================================
%% Closing
%% ===================================================================
%%
%%    \begin{macrocode}
%<!CLASS>
\ExplSyntaxOff
%    \end{macrocode}
%%
%% \Finale
%% \endinput
